% p.57 — الباب الثالث والعشرون في الناقص والأجوف في حالة الإعراب
% Layout (as in the lithograph): two tables printed SIDE BY SIDE under the chapter title in
% one `tables` row, each with its own caption row spanning its three columns. RIGHT half:
% «افعال ناقصة» (مرفوع، منصوب، مجزوم), with a second block headed «امر | نهى | امر غائب».
% LEFT half: «افعال اجوف» (مرفوع، اصل في الجزم، موجود), then a block «امر في الاصل | مبدل | موجود»,
% then a block «مثال ماضي | مثال مضارع | مثال موجود».
\mspage{57}
\bab{الباب الثالث والعشرون في الناقص والأجوف في حالة الإعراب}
\begin{tables}
\begin{sarf}{|c|c|c|}
\multicolumn{3}{|c|}{\thead{افعال ناقصة}}\\\hline
\thead{مَرْفُوْعٌ} & \thead{مَنْصُوْبٌ} & \thead{مَجْزُوْمٌ}\\\hline
يَضْرِبُ & لَنْ يَضْرِبَ & لَمْ يَضْرِبْ\\
يَضْرِبَانِ & لَنْ يَضْرِبَا & لَمْ يَضْرِبَا\\
يَضْرِبُوْنَ & لَنْ يَضْرِبُوْا & لَمْ يَضْرِبُوْا\\
يَضْرِبْنَ & لَنْ يَضْرِبْنَ & لَمْ يَضْرِبْنَ\\
تَضْرِبْنَ & لَنْ تَضْرِبْنَ & لَمْ تَضْرِبْنَ\\\hline
يَدْعُوْ & لَنْ يَدْعُوَ & لَمْ يَدْعُ\\
يَدْعُوَانِ & لَنْ يَدْعُوَا & لَمْ يَدْعُوَا\\
يَدْعُوْنَ & لَنْ يَدْعُوْا & لَمْ يَدْعُوْا\\
يَدْعُوْنَ & لَنْ يَدْعُوْنَ & لَمْ يَدْعُوْنَ\\
تَدْعُوْنَ & لَنْ تَدْعُوْنَ & لَمْ تَدْعُوْنَ\\
يَخْشٰى & لَنْ يَخْشٰى & لَمْ يَخْشَ\\
يَرْمِيْ & لَنْ يَرْمِيَ & لَمْ يَرْمِ\\\hline
\thead{امر} & \thead{نهى} & \thead{امر غائب}\\\hline
اِضْرِبْ & لَا تَضْرِبْ & لِيَضْرِبْ\\
اُدْعُ & لَا تَدْعُ & لِيَدْعُ\\
اُدْعُوْا & لَا تَدْعُوْا & لِيَدْعُوْا\\
اِخْشَ & لَا تَخْشَ & لِيَخْشَ\\
اِخْشَيَا & لَا تَخْشَيَا & لِيَخْشَيَا\\
اِرْمِ & لَا تَرْمِ & لِيَرْمِ\\
\end{sarf}
\begin{sarf}{|c|c|c|}
\multicolumn{3}{|c|}{\thead{افعال اجوف}}\\\hline
\thead{مَرْفُوْعٌ} & \thead{اَصْلٌ فِي الْجَزْمِ} & \thead{مَوْجُوْدٌ}\\\hline
يَقُوْلُ & لَمْ يَقُوْلْ & لَمْ يَقُلْ\\
يَقُوْلَانِ & لَمْ يَقُوْلَا & لَمْ يَقُوْلَا\\
يَقُلْنَ & لَمْ يَقُوْلْنَ & لَمْ يَقُلْنَ\\
يَخَافُ & لَمْ يَخَافْ & لَمْ يَخَفْ\\
يَخَافَانِ & لَمْ يَخَافَا & لَمْ يَخَافَا\\
يَخَفْنَ & لَمْ يَخَافْنَ & لَمْ يَخَفْنَ\\
يَبِيْعُ & لَمْ يَبِيْعْ & لَمْ يَبِعْ\\
يَبِيْعَانِ & لَمْ يَبِيْعَا & لَمْ يَبِيْعَا\\
يَبِعْنَ & لَمْ يَبِيْعْنَ & لَمْ يَبِعْنَ\\\hline
\thead{امر في الاصل} & \thead{مبدل} & \thead{موجود}\\\hline
اُقْوُلْ & قُوْلْ & قُلْ\\
اُقْوُلَا & قُوْلَا & قُوْلَا\\
اُقْوُلْنَ & قُوْلْنَ & قُلْنَ\\
اِبْيِعْ & بِيْعْ & بِعْ\\\hline
\thead{مثال ماضي} & \thead{مثال مضارع} & \thead{مثال موجود}\\\hline
وَقَعَ & يَوْقَعُ & يَقَعُ\\
وَعَدَ & يَوْعِدُ & يَعِدُ\\
وَرِمَ & يَوْرِمُ & يَرِمُ\\
وَجِلَ & يَوْجَلُ & يَوْجَلُ\\
\end{sarf}
\end{tables}
